\documentclass[12pt]{article}
\usepackage[T1]{fontenc}
\usepackage[utf8]{inputenc}
\usepackage{lmodern}
\usepackage{textcomp}
\usepackage{enumitem}
\usepackage{geometry}
\geometry{margin=1in}
\begin{document}
\begin{center}
Answer Key - Constitutional Law - Undergraduate Level Assessment 9 \
\small (Answers are ordered exactly as in the question paper.)
\end{center}

\section*{Multiple Choice}
\begin{enumerate}[label=\arabic*.]
\item \textbf{Answer:} C
\\ \textit{Options:} A) Domestic law always prevails over international law; B) Only customary international law prevails over domestic law; C) Obligations under international law prevail over domestic law; D) Constitutional obligations always prevail over obligations under international law
\\ \textit{Note:} Correct option: C - Obligations under international law prevail over domestic law
\item \textbf{Answer:} D
\\ \textit{Options:} A) The ultimate source of a legal system's morality.; B) The rule that distinguishes norms from habits of obedience.; C) The constitution of a state.; D) A presupposition that facilitates our understanding of the legal system.
\\ \textit{Note:} Correct option: D - A presupposition that facilitates our understanding of the legal system.
\item \textbf{Answer:} B
\\ \textit{Options:} A) Piracy jure gentium is subject to flag State jurisdiction; B) Piracy jure gentium is subject to universal jurisdiction; C) Piracy jure gentium is subject to port State jurisdiction; D) Piracy jure gentium is subject to nationality-based jurisdiction
\\ \textit{Note:} Correct option: B - Piracy jure gentium is subject to universal jurisdiction
\item \textbf{Answer:} A
\\ \textit{Options:} A) He denies that law develops independently of social and economic forces.; B) He claims that law is economically immoral.; C) He rejects a positivist account of law.; D) He opposes any sociological analysis of the law.
\\ \textit{Note:} Correct option: A - He denies that law develops independently of social and economic forces.
\item \textbf{Answer:} B
\\ \textit{Options:} A) John Rawls; B) Stammler; C) Jerome Hall; D) John Finns
\\ \textit{Note:} Correct option: B - Stammler
\end{enumerate}

\section*{True / False}
\begin{enumerate}[label=\arabic*.]
\setcounter{enumi}{5}
\item \textbf{Answer:} False
\\ \textit{Options:} A) True; B) False
\\ \textit{Note:} LLM-generated false statement. Correct answer: 'Gemeinschaft is a society based on community, Gesellschaft one based on association'.
\item \textbf{Answer:} True
\\ \textit{Options:} A) True; B) False
\\ \textit{Note:} LLM-generated true statement based on MCQ.
\item \textbf{Answer:} False
\\ \textit{Options:} A) True; B) False
\\ \textit{Note:} LLM-generated false statement. Correct answer: 'Salmond'.
\item \textbf{Answer:} True
\\ \textit{Options:} A) True; B) False
\\ \textit{Note:} LLM-generated true statement based on MCQ.
\item \textbf{Answer:} True
\\ \textit{Options:} A) True; B) False
\\ \textit{Note:} LLM-generated true statement based on MCQ.
\end{enumerate}

\section*{Long Form Response}
\begin{enumerate}[label=\arabic*.]
\setcounter{enumi}{10}
\item \textbf{Answer:} The UK Constitution is uncodified and can be found in a number of sources. Explain why this is the correct answer and provide supporting evidence. Reference key facts or concepts that support this choice.
\\ \textit{Note:} Long-form derived from MCQ; award credit for stating the correct idea and giving supporting reasoning.
\item \textbf{Answer:} Mechanical solidarity leads to repressive law. Organic solidarity leads to restitutive law.. Explain why this is the correct answer and provide supporting evidence. Reference key facts or concepts that support this choice.
\\ \textit{Note:} Long-form derived from MCQ; award credit for stating the correct idea and giving supporting reasoning.
\end{enumerate}

\vfill
\noindent\textit{End of Answer Key}
\end{document}